量子算法处理的是可控制的概率振幅, 而不是把所有经典答案同时读出. 量子信息的任务是说明哪些相干与纠缠资源能提高计算、通信或保护信息的能力, 并把测量和噪声也纳入同一数学模型. 读者应熟悉线性代数、量子态、算符、复合系统与基本测量; 这些先修见本仓库量子力学卷. 本卷从量子电路的计算语义开始, 沿算法、信道、编码到多体资源与通信, 在新增定义处给可复算的例子.

\textbf{课程与联列去重.} 核心版本为Peter W. Shor的\href{https://ocw.mit.edu/courses/18-435j-quantum-computation-fall-2003/}{18.435J / 2.111, Fall 2003}及Isaac Chuang、Peter Shor的\href{https://ocw.mit.edu/courses/mas-865j-quantum-information-science-spring-2006/}{MAS.865J, Spring 2006}. 后者原件联列6.443J / 8.371J / 18.409 / MAS.865, 只计一门同课, 不以四编号重复统计. 2006的8.371J与2018的8.371x不能仅凭编号视为同一份资料. 2018的8.370x与8.371x的OCW入口和Open Learning Library公开入口已核查, 其网页并无可直接等同本卷题库的本地作业PDF; 不把课程介绍页计为已译整门课程.

\textbf{资料层级.} 18.435J公开的是部分课笔记和考核, 不是23次课各有完整笔记. MAS.865J有20次理论课学生笔记, 后六次为项目会面或展示, 不是缺失理论章. 笔记者为Kayla Jacobs、Jennifer Novosad、Saikat Guha、Igor Sylvester(原件有Sylvestor拼写)、Danial Lashkari、Salman Abolfathe、Edward Platt、Isaac Kim、Stephen Jordan. 笔记由OCW经许可发布, 仍可能含排印与推导错误; 本卷在使用处独立复核, 不以“官方托管”代替数学检查.

\textbf{题目与解答身份.} 官方公开题面译写保留所有可用条件. 有官方解答的题目注明参考官方解答重编, 更正与补推导另标; 其余为编者解答. 商业教材只有题号引用或直接转载的段落不自动取得本项目再出版许可, 只留书目定位与缺项, 不补造题面. 已撤回的问题保持撤回, 开放研究问题不虚称已找到最优. 数值与电路检查支持具体实例, 不替代一般证明.

\textbf{许可与署名.} OCW及以上原作者保留权利. 对可改编OCW内容, 本卷翻译、重排、推导补写与自绘图采用\href{https://creativecommons.org/licenses/by-nc-sa/4.0/}{CC BY-NC-SA 4.0}, 非商业使用、保留署名并相同方式共享. 单项第三方例外不由默认许可覆盖. 本项目不表示MIT、教师或笔记者认可. 逐文件URL与哈希见来源清单, 使用条款见\url{https://ocw.mit.edu/pages/privacy-and-terms-of-use/}.
